% =============================================================================
% HiddenPower — clean XeLaTeX preamble
% =============================================================================

% -----------------------------------------------------------------------------
% Language and fonts
% -----------------------------------------------------------------------------

\usepackage{fontspec}

\defaultfontfeatures{
    Ligatures=TeX,
    Scale=MatchLowercase
}

\IfFontExistsTF{CMU Serif}{
    \setmainfont{CMU Serif}
}{
    \setmainfont{DejaVu Serif}
}

\IfFontExistsTF{CMU Sans Serif}{
    \setsansfont{CMU Sans Serif}
}{
    \setsansfont{DejaVu Sans}
}

\IfFontExistsTF{CMU Typewriter Text}{
    \setmonofont{CMU Typewriter Text}
}{
    \setmonofont{DejaVu Sans Mono}
}

\usepackage[english,russian]{babel}
\usepackage{microtype}

% -----------------------------------------------------------------------------
% Page geometry and typography
% -----------------------------------------------------------------------------

\usepackage[
    a4paper,
    inner=30mm,
    outer=25mm,
    top=25mm,
    bottom=28mm,
    headheight=16pt,
    headsep=8mm,
    footskip=12mm
]{geometry}

\usepackage{setspace}
\setstretch{1.05}

\setlength{\parindent}{1.5em}
\setlength{\parskip}{0pt}
\setlength{\emergencystretch}{2em}

\raggedbottom

\setcounter{secnumdepth}{3}
\setcounter{tocdepth}{1}

% -----------------------------------------------------------------------------
% Core mathematics
% -----------------------------------------------------------------------------

\usepackage{amsmath}
\usepackage{amssymb}
\usepackage{amsthm}
\usepackage{mathtools}

\usepackage{bm}
\usepackage{mathrsfs}
\usepackage{esint}

\allowdisplaybreaks[2]
\setlength{\jot}{6pt}
\numberwithin{equation}{section}

% -----------------------------------------------------------------------------
% Figures and tables
% -----------------------------------------------------------------------------

\usepackage{graphicx}
\graphicspath{{figures/}}

\usepackage{float}
\usepackage{caption}
\usepackage{subcaption}

\usepackage{booktabs}
\usepackage{array}
\usepackage{tabularx}

% -----------------------------------------------------------------------------
% Source code
% -----------------------------------------------------------------------------

\usepackage{xcolor}
\usepackage{listings}

\lstdefinestyle{hiddenpower}{
    basicstyle=\ttfamily\small,
    numbers=left,
    numberstyle=\scriptsize,
    numbersep=8pt,
    frame=single,
    breaklines=true,
    showstringspaces=false,
    tabsize=4,
    columns=fullflexible
}

\lstset{style=hiddenpower}

% -----------------------------------------------------------------------------
% Full-page cover image
% -----------------------------------------------------------------------------

\usepackage{eso-pic}

\newcommand{\fullpageimage}[1]{%
    \clearpage
    \thispagestyle{empty}%
    \AddToShipoutPictureBG*{%
        \AtPageLowerLeft{%
            \includegraphics[
                width=\paperwidth,
                height=\paperheight
            ]{#1}%
        }%
    }%
    \null
    \clearpage
}

% -----------------------------------------------------------------------------
% Theorem system
% -----------------------------------------------------------------------------

\newtheoremstyle{strictplain}
    {8pt}
    {8pt}
    {\itshape}
    {}
    {\bfseries}
    {.}
    {0.5em}
    {}

\newtheoremstyle{strictdefinition}
    {8pt}
    {8pt}
    {\normalfont}
    {}
    {\bfseries}
    {.}
    {0.5em}
    {}

\newtheoremstyle{strictremark}
    {6pt}
    {6pt}
    {\normalfont}
    {}
    {\itshape}
    {.}
    {0.5em}
    {}

\theoremstyle{strictplain}

\newtheorem{theorem}{Теорема}[section]
\newtheorem{lemma}[theorem]{Лемма}
\newtheorem{proposition}[theorem]{Предложение}
\newtheorem{corollary}[theorem]{Следствие}
\newtheorem{criterion}[theorem]{Критерий}
\newtheorem{claim}[theorem]{Утверждение}

\theoremstyle{strictdefinition}

\newtheorem{definition}[theorem]{Определение}
\newtheorem{assumption}[theorem]{Предположение}
\newtheorem{construction}[theorem]{Конструкция}
\newtheorem{principle}[theorem]{Принцип}
\newtheorem{example}[theorem]{Пример}
\newtheorem{counterexample}[theorem]{Контрпример}
\newtheorem{problem}[theorem]{Задача}
\newtheorem{exercise}[theorem]{Упражнение}

\theoremstyle{strictremark}

\newtheorem{remark}[theorem]{Замечание}
\newtheorem{notation}[theorem]{Обозначение}
\newtheorem{convention}[theorem]{Соглашение}
\newtheorem{caution}[theorem]{Предупреждение}

\newtheorem{connection}[theorem]{Связь с другими разделами}
\newtheorem{mlconnection}[theorem]{Связь с машинным обучением}
\newtheorem{application}[theorem]{Приложение}
\newtheorem{computationalnote}[theorem]{Вычислительная реализация}

\newenvironment{derivation}
    {\begin{proof}[Вывод]}
    {\end{proof}}

\newenvironment{proofidea}
    {\begin{proof}[Идея доказательства]}
    {\end{proof}}

\renewcommand{\proofname}{Доказательство}
\renewcommand{\qedsymbol}{\(\square\)}

% -----------------------------------------------------------------------------
% Safe mathematical notation
% -----------------------------------------------------------------------------

\providecommand{\N}{\mathbb{N}}
\providecommand{\Z}{\mathbb{Z}}
\providecommand{\Q}{\mathbb{Q}}
\providecommand{\R}{\mathbb{R}}
\providecommand{\C}{\mathbb{C}}
\providecommand{\F}{\mathbb{F}}

\providecommand{\vct}[1]{\bm{#1}}
\providecommand{\mat}[1]{\bm{#1}}
\providecommand{\ten}[1]{\bm{\mathcal{#1}}}

\providecommand{\trans}{\mathsf{T}}
\providecommand{\adj}{\mathsf{H}}

\providecommand{\dd}{\mathop{}\!\mathrm{d}}
\providecommand{\given}{\,\mid\,}

\providecommand{\mathbbP}{\mathbb{P}}
\providecommand{\mathbbE}{\mathbb{E}}

\providecommand{\rankop}{\operatorname{rank}}
\providecommand{\traceop}{\operatorname{tr}}
\providecommand{\diagop}{\operatorname{diag}}
\providecommand{\spanop}{\operatorname{span}}
\providecommand{\projop}{\operatorname{proj}}
\providecommand{\kerop}{\operatorname{Ker}}
\providecommand{\imop}{\operatorname{Im}}

\providecommand{\argmin}{\operatorname*{arg\,min}}
\providecommand{\argmax}{\operatorname*{arg\,max}}

\providecommand{\grad}{\nabla}
\providecommand{\laplacian}{\Delta}

\providecommand{\bigO}{\mathcal{O}}
\providecommand{\littleo}{o}

% -----------------------------------------------------------------------------
% Delimiters
% -----------------------------------------------------------------------------

\DeclarePairedDelimiter{\abs}{\lvert}{\rvert}
\DeclarePairedDelimiter{\norm}{\lVert}{\rVert}
\DeclarePairedDelimiter{\paren}{\lparen}{\rparen}
\DeclarePairedDelimiter{\bracket}{\lbrack}{\rbrack}
\DeclarePairedDelimiter{\floor}{\lfloor}{\rfloor}
\DeclarePairedDelimiter{\ceil}{\lceil}{\rceil}

\DeclarePairedDelimiterX{\inner}[2]
    {\langle}{\rangle}
    {#1,#2}

\DeclarePairedDelimiterX{\setbuild}[2]
    {\lbrace}{\rbrace}
    {#1\,\middle|\,#2}

% -----------------------------------------------------------------------------
% Differential notation
% -----------------------------------------------------------------------------

\newcommand{\dv}[2]{%
    \frac{\dd #1}{\dd #2}
}

\newcommand{\pdv}[2]{%
    \frac{\partial #1}{\partial #2}
}

\newcommand{\fdv}[2]{%
    \frac{\delta #1}{\delta #2}
}

\newcommand{\Jac}[1]{%
    \mathrm{J}_{#1}
}

\newcommand{\Hessian}[1]{%
    \nabla^{2}#1
}

% -----------------------------------------------------------------------------
% Headers
% -----------------------------------------------------------------------------

\usepackage{fancyhdr}
\usepackage{emptypage}

\pagestyle{fancy}
\fancyhf{}

\fancyhead[L]{%
    \nouppercase{\leftmark}
}

\fancyhead[R]{%
    \thepage
}

\renewcommand{\headrulewidth}{0.4pt}
\renewcommand{\footrulewidth}{0pt}

\renewcommand{\chaptermark}[1]{%
    \markboth{\thechapter.\ #1}{}
}

\renewcommand{\sectionmark}[1]{%
    \markright{\thesection.\ #1}
}

% -----------------------------------------------------------------------------
% Hyperlinks and references
% Keep this order: hyperref -> bookmark -> cleveref
% -----------------------------------------------------------------------------

\usepackage{xurl}

\usepackage[
    unicode=true,
    colorlinks=true,
    linkcolor=black,
    citecolor=black,
    urlcolor=blue,
    pdfborder={0 0 0}
]{hyperref}

\usepackage{bookmark}

\usepackage[
    nameinlink,
    capitalise,
    noabbrev
]{cleveref}

\crefname{theorem}{теорема}{теоремы}
\crefname{lemma}{лемма}{леммы}
\crefname{proposition}{предложение}{предложения}
\crefname{corollary}{следствие}{следствия}
\crefname{criterion}{критерий}{критерии}
\crefname{claim}{утверждение}{утверждения}

\crefname{definition}{определение}{определения}
\crefname{assumption}{предположение}{предположения}
\crefname{construction}{конструкция}{конструкции}
\crefname{principle}{принцип}{принципы}
\crefname{example}{пример}{примеры}
\crefname{counterexample}{контрпример}{контрпримеры}
\crefname{problem}{задача}{задачи}
\crefname{exercise}{упражнение}{упражнения}

\crefname{remark}{замечание}{замечания}
\crefname{notation}{обозначение}{обозначения}
\crefname{convention}{соглашение}{соглашения}
\crefname{caution}{предупреждение}{предупреждения}

\crefname{connection}{связь}{связи}
\crefname{mlconnection}
    {связь с машинным обучением}
    {связи с машинным обучением}
\crefname{application}{приложение}{приложения}
\crefname{computationalnote}
    {вычислительная реализация}
    {вычислительные реализации}

\crefname{equation}{формула}{формулы}
\crefname{figure}{рисунок}{рисунки}
\crefname{table}{таблица}{таблицы}

% BEGIN HIDDENPOWER TOC PAGE NUMBER WIDTH
% Стандартного поля 1.55em недостаточно для трёхзначных номеров страниц
% при основном шрифте DejaVu Serif.
\makeatletter
\renewcommand{\@pnumwidth}{2.2em}
\renewcommand{\@tocrmarg}{3.2em}
\makeatother
% END HIDDENPOWER TOC PAGE NUMBER WIDTH

% -----------------------------------------------------------------------------
% Page-break discipline
% -----------------------------------------------------------------------------

\clubpenalty=10000
\widowpenalty=10000
\displaywidowpenalty=10000
\interfootnotelinepenalty=10000

% BEGIN HIDDENPOWER TOC NUMBER WIDTH FIX
% Стандартные поля book.cls слишком узки для номеров вида 1.10 и 1.10.10.
% Начало текста каждого следующего уровня выровнено по началу текста
% предыдущего уровня.
\makeatletter
\renewcommand*\l@section{%
    \@dottedtocline{1}{1.5em}{3.6em}%
}
\renewcommand*\l@subsection{%
    \@dottedtocline{2}{5.1em}{4.5em}%
}
\renewcommand*\l@subsubsection{%
    \@dottedtocline{3}{9.6em}{5.2em}%
}
\makeatother
% END HIDDENPOWER TOC NUMBER WIDTH FIX

% BEGIN HIDDENPOWER LIGHT TOC CLEANUP
% В оглавлении показываются только главы и разделы. Малый кегль и локальная
% растяжимость строк устраняют Overfull в длинных математических заголовках.
\makeatletter
\let\hiddenpoweroriginaltableofcontents\tableofcontents
\renewcommand{\tableofcontents}{%
    \begingroup
    \small
    \setlength{\emergencystretch}{3em}%
    \sloppy
    \hiddenpoweroriginaltableofcontents
    \endgroup
}
\makeatother
% END HIDDENPOWER LIGHT TOC CLEANUP
